\chapter{Conclusions and future developments}
\label{ch:conclusions}%
% Aim: 4-6 pages.

\section{Main findings}
% Answer the objectives stated in the Introduction, one by one, with numbers.
% Draft, to be revised when the test campaign is complete.
\begin{itemize}
    \item Under constant oedometric stress, dry LHS-1 compacts only during loading, while mixtures of LHS-1 and ice keep compacting for days: the creep is controlled by the ice.
    \item About 10\% of ice by mass is enough to produce creep ($C_\alpha\approx0.04$ at $-20\,^\circ$C and 2\,MPa); at equal void ratio the rate increases by orders of magnitude as the ice fraction rises towards pore filling.
    \item The early compaction is independent of temperature, while the creep is thermally activated, with a strong increase between $-20$ and $-10\,^\circ$C.
    \item Because the oedometric strain is entirely volumetric, the isochoric flow laws of ice cannot describe the tests on their own; a compressible porous law with an ice matrix following Goldsby and Kohlstedt is adopted.
\end{itemize}

\section{Limitations}

\section{Future developments}
% Triaxial tests with controlled radial stress; forced-oscillation (PinQ) measurements of elastic moduli;
% lower temperatures; other simulants (LMS-1) and ice morphologies; implementation in a finite-element code.
